% Generated draft: requires mathematical and source review.
\documentclass[11pt]{article}
\usepackage[margin=25mm]{geometry}
\usepackage{fontspec}
\setmainfont{Latin Modern Roman}
\usepackage{amsmath,amssymb,mathtools,unicode-math}
\setmathfont{Latin Modern Math}
\usepackage{enumitem,needspace,xurl,hyperref}
\setlength{\parindent}{0pt}
\setlength{\parskip}{5pt}
\setlength{\emergencystretch}{3em}
\newcommand{\N}{\mathbb{N}}
\newcommand{\Z}{\mathbb{Z}}
\newcommand{\Q}{\mathbb{Q}}
\newcommand{\R}{\mathbb{R}}
\newcommand{\C}{\mathbb{C}}
\newcommand{\F}{\mathbb{F}}
\newcommand{\A}{\mathbb{A}}
\newcommand{\PP}{\mathbb{P}}
\newcommand{\E}{\mathbb{E}}
\newcommand{\eps}{\varepsilon}
\newcommand{\dd}{\,\mathrm{d}}
\begin{document}
\section*{UnsolvedMath ID: 20001066}
\textbf{Source problem code:} catalogue entry
\subsection{Definitions and mathematical statement}
The source's Conjecture 6.21, attributed to Hamidoune (1981), states: if a digraph $D$ has minimum outdegree and indegree at least $r\ge1$, then some arc $x\to y$ has reverse local vertex-connectivity $\kappa_D(y,x)\ge r$. Here $\kappa_D(y,x)$ is the maximum number of internally vertex-disjoint directed $y$-to-$x$ paths. The original immediately notes Thomassen's Theorem 5.2 as a counterexample for every $r\ge3$, so its conjecture wording is historical rather than an unresolved assertion in those cases.

The attached aid emphasizes the source's permitted parallel arcs. Under finite loopless multidigraph conventions with degrees counting arc multiplicity, for every $m\ge3$ and $r\ge2$ take a directed $m$-cycle and replace each arc by $r$ parallel copies. Every vertex has arc indegree and outdegree $r$, but for every selected arc $x\to y$, all reverse paths follow the unique support route and share an internal vertex. Thus
\[
\kappa_D(y,x)=1<r
\]
for every arc. The example is loopless and avoids any ambiguity about counting parallel direct reverse paths, since the reverse route has length at least two.

For $r=1$ in the finite loopless setting, a directed cycle supplies an arc with a reverse path, so the assertion holds. For a repair define support outdegree and indegree by numbers of distinct out- and in-neighbors. The aid states that the universal support-degree version is logically equivalent to the simple-digraph version: a simple support satisfying the conjecture supplies paths that lift to any parallel-arc realization, and simple digraphs are included among multidigraphs. This equivalence does not require claiming that every individual connectivity value is unchanged by parallel arcs.

For the repaired $r=2$ statement the aid gives the bidirected-support case, where each support adjacency is present in both directions and every vertex has at least two distinct neighbors. It records the general repaired simple-digraph $r=2$ case as unresolved in the literature checked, without a full citation. Hence the literal multiplicity-degree refutation must not be presented as a refutation of that separate repaired case.
\subsection{Short English statement}
The reverse-connectivity assertion depends on degree conventions. The supplied aid gives parallel-cycle counterexamples for every $r\ge2$ under multiplicity degrees, while separating the true $r=1$ case, the known historical $r\ge3$ refutation and the distinct support-degree repair.
\subsection{Sources}
\begin{itemize}
\item Local UnsolvedMath registry, permanent ID $20001066$, source code catalogue entry; canonical \texttt{aim-combinatorics-notes.json}, \texttt{notes[190]}.
\item American Institute of Mathematics, ``The Caccetta--H\"aggkvist conjecture'' workshop record, Conjecture 6.21 attributed to Hamidoune (1981): \url{https://aimath.org/WWN/caccetta/caccetta.pdf}. Its own note cites Thomassen, Theorem 5.2, for counterexamples at all $r\ge3$, without the complete bibliographic entry here. No PDF inspection is claimed.
\item Attached research attempt 1: counterexample, recorded run status refuted, rigor proved, novelty confidence low; registry status solved, classification SOLVED-IN-LITERATURE. Its convention analysis and repaired-case assessment are machine generated and not peer reviewed; independently verify before citation.
\end{itemize}
\end{document}
